\section{Elementary finite arithmetic}
\label{sec:finite}

We work with ordinary finite sets, sums and products as readable notation
for bounded loops. Empty sums are \(0\), empty products are \(1\), and
squarefreeness always refers to a \emph{positive} integer.
All rational denominators used below are proved positive.

\subsection{Division, prime factors, and finite coding}

For \(d>0\), division with remainder is obtained by a loop on the dividend.
Starting from quotient and remainder \((0,0)\), increase the remainder by
one at each step; when it reaches \(d\), reset it to \(0\) and increase
the quotient. Induction proves
\[
 n=dq+r,\qquad 0\le r<d.
\]
The stored pair can be coded by a primitive recursive pairing function.
No unbounded search is involved.

For \(n>1\), search the finite list \(2,\ldots,n\) for its first divisor.
The search succeeds because \(n\) is on the list. Its first result \(p\)
is prime: a divisor \(a\) of \(p\) with \(1<a<p\) would also divide \(n\)
and would have been found first. Divide out \(p\) and repeat, leaving the
state unchanged once the remaining factor is \(1\). Every successful step
strictly reduces that positive remaining factor, so a loop of \(n\) steps
suffices. Its invariant is
\[
 n=(\hbox{current remaining factor})
       \prod(\hbox{prime factors already recorded}).
\]
This constructs a prime-factor list, including multiplicities, by primitive
recursion.

For completeness, the divisibility fact used to prove uniqueness is
constructive. The subtraction version of the Euclidean algorithm maintains
the common divisors of a pair of nonnegative integers; whenever both
coordinates are positive, subtract the smaller from the larger. Their sum
strictly decreases, so a fixed loop of \(a+b\) steps, frozen at termination,
computes \(\gcd(a,b)\). Maintaining integer linear-combination coefficients
at the same time gives a Bezout identity. If \(p\) is prime and \(p\nmid a\),
then \(\gcd(p,a)=1\), so some explicitly computed integers \(u,v\) satisfy
\(up+va=1\). Multiplying by \(b\) shows
\[
 p\mid ab\quad\Longrightarrow\quad p\mid a\ \hbox{or}\ p\mid b.
\]
Here the case split \(p\mid a\) is a decidable divisibility test. Repeatedly
remove a matching prime from two factor lists to obtain uniqueness.
All the lists and coefficient traces have fixed finite bounds; a more
explicit coding audit is given in Section~\ref{sec:ipra}.

We may consequently use:
\[
 \rad(n)=\prod_{p\mid n}p,\qquad
 \omega(n)=\#\{p:p\mid n\}
 \quad(n\ge1),
\]
with \(\rad(1)=1\), and the Mobius function
\[
 \mu(n)=
 \begin{cases}
 (-1)^{\omega(n)},&n\text{ squarefree},\\
 0,&n\text{ not squarefree}.
 \end{cases}
\]
These are all obtained from bounded factorization. In particular,
\(\rad(n)\mid n\), and \(1\le\rad(n)\le n\).

\subsection{The residue counts}

For squarefree \(d\), define
\[
 \rho(d)=\#\{a:0\le a<d,\ d\mid a(a+2)\}.
\]
For a prime \(p\), the preceding divisibility fact gives the two possible
root residues \(0\) and \(-2\). They coincide exactly when \(p=2\).
Thus
\begin{equation}
 \rho(2)=1,\qquad
 \rho(p)=2\quad(p\text{ odd prime}),\qquad
 \rho(1)=1.
 \label{eq:local-rho}
\end{equation}

We recall the finite Chinese remainder bijection rather than treating it
as an existence principle. If \(a,b>0\) are coprime, compute an inverse
\(u\) of \(a\) modulo \(b\) from a Bezout identity. Given residues
\(0\le r<a\), \(0\le s<b\), set
\[
 x=r+a\bigl((s-r)u\bmod b\bigr).
\]
The remainder is read as its representative in \(\{0,\ldots,b-1\}\);
signed integers can be coded by pairs of natural numbers.
Then \(0\le x<ab\), \(x\equiv r\pmod a\), and \(x\equiv s\pmod b\).
If two representatives have these properties, their difference is
divisible by \(a\) and \(b\), hence by \(ab\): if the difference is \(ak\),
then \(b\mid ak\), and a Bezout identity for \(a,b\) implies \(b\mid k\).
Its absolute value is below \(ab\), so it is zero. This proves the bijection.

Applying the bijection successively to the distinct prime factors of \(d\)
gives
\begin{equation}
 \rho(d)=\prod_{p\mid d}\rho(p).
 \label{eq:rho-product}
\end{equation}
For squarefree \(d,e\), the local factors are at least \(1\), so
\begin{equation}
 \rho(\lcm(d,e))\le\rho(d)\rho(e).
 \label{eq:rho-lcm}
\end{equation}

\begin{lemma}[Finite counting error]
\label{lem:count-error}
For \(X\ge0\) and squarefree \(d\), let
\[
 A_d(X)=\#\{n:1\le n\le X,\ d\mid n(n+2)\}.
\]
Then
\begin{equation}
 \left|A_d(X)-X\frac{\rho(d)}d\right|\le\rho(d).
 \label{eq:count-error}
\end{equation}
\end{lemma}

\begin{proof}
Write \(X=qd+r\), \(0\le r<d\). The interval \(1,\ldots,X\)
consists of \(q\) complete consecutive blocks of \(d\) integers and a
remaining block of \(r\) integers. Each complete block contains exactly
one representative of each residue. A fixed root residue therefore occurs
either \(q\) or \(q+1\) times. Its difference from \(X/d=q+r/d\)
has absolute value at most \(1\). Sum this estimate over the
\(\rho(d)\) root residues and use the finite triangle inequality.
For \(d=1\) the error is exactly \(0\).
\end{proof}

\subsection{The cancellation identity}

\begin{lemma}[Squarefree Mobius cancellation]
\label{lem:mobius}
For squarefree \(v\),
\[
 \sum_{a\mid v}\mu(a)=
 \begin{cases}1,&v=1,\\0,&v>1.\end{cases}
\]
If \(r\mid v\) and \(v\) is squarefree, then
\begin{equation}
 \sum_{r\mid d\mid v}\mu(d)=
 \begin{cases}\mu(r),&v=r,\\0,&v\ne r.\end{cases}
 \label{eq:interval-mobius}
\end{equation}
\end{lemma}

\begin{proof}
Divisors of \(v\) correspond to subsets of its finite list of prime factors.
Expanding \(\prod_{p\mid v}(1-1)\) gives the first sum. For the second,
write \(d=ra\). Since \(v\) is squarefree, \(r\) and \(v/r\) are coprime,
so \(\mu(ra)=\mu(r)\mu(a)\) for \(a\mid v/r\). Apply the first identity.
The product expansion can also be proved directly by induction on the
length of the prime-factor list.
\end{proof}
